% Seção elaborada a partir de fontes públicas consultadas em 25/10/2026.
% NÃO representa testes com usuários, telemetria ou acesso interno às plataformas.
\section{Benchmarking e Matriz AARR}
\label{sec:benchmarking}

\subsection{Objetivo e procedimento de análise}
O \textit{benchmarking} consiste em comparar características de soluções de referência para identificar práticas, alternativas de projeto e oportunidades de melhoria. No projeto \textit{Aula em Dia}, o objetivo é observar como ferramentas existentes atendem às atividades de elaboração, organização e reutilização de planos de aula, além de reconhecer possibilidades de uso responsável de inteligência artificial no planejamento docente.

Foi conduzida uma \textbf{análise documental exploratória} de informações públicas, disponibilizadas nas páginas oficiais de três plataformas: Teachy, MagicSchool AI e Planbook. A seleção buscou contemplar uma solução brasileira com geração de materiais pedagógicos por IA (Teachy), uma solução especializada em ferramentas de IA para professores (MagicSchool) e uma ferramenta com ênfase na gestão recorrente de planos de aula (Planbook). Os recursos foram registrados com base em páginas consultadas em \textbf{25 de outubro de 2026}. Não foram realizados testes de usabilidade, avaliações com professores nem medições de desempenho ou de adoção; por isso, as observações representam funcionalidades \emph{documentadas pelos fornecedores}, e não eficácia comprovada.

\subsection{Descrição das plataformas de referência}
\begin{description}[leftmargin=!,labelwidth=3cm]
  \item[Teachy.] Plataforma voltada a professores, cuja documentação apresenta planejamento por período, associação a turmas e disciplinas, calendário, edição dos planejamentos gerados e criação de planos de aula com apoio de IA. As orientações publicadas também mencionam objetivos associados à Base Nacional Comum Curricular (BNCC) \citep{teachy_planejamento2026,teachy_ferramentas2026}.
  \item[MagicSchool AI.] Plataforma educacional que disponibiliza ferramentas de geração assistida de planos de aula, atividades e avaliações. Sua ferramenta de planejamento permite informar disciplina, etapa escolar, objetivos e referenciais curriculares e, posteriormente, revisar o material produzido \citep{magicschool_planos2026,magicschool_ferramentas2026}.
  \item[Planbook.] Solução de planejamento de aulas com organização de calendário, reutilização de planos e modelos, anexação de recursos, referências curriculares e compartilhamento controlado com outros professores \citep{planbook_professores2026}.
\end{description}

\subsection{Matriz comparativa das funcionalidades}
A Tabela~\ref{tab:comparacao-benchmark} sintetiza apenas características observadas na documentação pública. A expressão ``não avaliado'' significa que não houve comprovação suficiente no conjunto de páginas consultadas; ela \textbf{não equivale} à inexistência da funcionalidade.

\begin{table}[htbp]
\centering
\caption{Comparação documental das soluções de referência}
\label{tab:comparacao-benchmark}
\footnotesize
\renewcommand{\arraystretch}{1.3}
\begin{tabularx}{\textwidth}{@{}p{2.8cm}XXX@{}}
\toprule
\textbf{Critério} & \textbf{Teachy} & \textbf{MagicSchool AI} & \textbf{Planbook} \\
\midrule
Elaboração de planos de aula & Sim, integrada ao planejamento & Sim, com ferramenta de IA & Sim, por editor e modelos \\
Geração de rascunhos com IA & Documentada & Documentada & Não avaliada \\
Organização temporal e calendário & Documentada & Não avaliada & Documentada \\
Revisão e adaptação pelo professor & Documentada & Documentada & Documentada \\
Referência explícita à BNCC brasileira & Documentada & Não avaliada & Não avaliada \\
Compartilhamento ou colaboração & Planejamento colaborativo anunciado & Coleções compartilháveis de ferramentas & Compartilhamento de planos e turmas \\
\bottomrule
\end{tabularx}
\par\smallskip
\begin{minipage}{\textwidth}\footnotesize Fonte: elaboração própria, com base na documentação oficial de \citet{teachy_planejamento2026,teachy_ferramentas2026,magicschool_planos2026,magicschool_ferramentas2026,planbook_professores2026}.\end{minipage}
\end{table}

\subsection{Matriz AARR: análise da jornada de uso}
O enunciado da disciplina solicita uma \emph{Matriz AARR}, mas não apresenta a definição de suas quatro letras. Para elaborar uma \textbf{versão preliminar explicitamente adaptada}, foram selecionadas quatro dimensões de jornada do usuário: \textbf{Aquisição}, \textbf{Ativação}, \textbf{Retenção} e \textbf{Recomendação}. A referência metodológica é o funil \textit{AARRR} (métricas piratas), cuja forma clássica compreende cinco dimensões, acrescentando \textbf{Receita} \citep{amazon_aarrr}. Neste relatório, a dimensão Receita não integra a matriz didática de quatro letras porque ainda não há modelo de monetização definido para o \textit{Aula em Dia}. \textbf{Essa interpretação precisa ser confirmada com o material da disciplina antes da entrega.}

A comparação AARR descrita na Tabela~\ref{tab:matriz-aarr} é qualitativa. Mostra elementos da experiência anunciados pelos serviços e oportunidades previstas para o projeto, \textbf{sem atribuir taxas de conversão, retenção ou recomendação} a plataformas externas.

\begin{table}[htbp]
\centering
\caption{Matriz AARR adaptada ao benchmarking do Aula em Dia}
\label{tab:matriz-aarr}
\scriptsize
\renewcommand{\arraystretch}{1.3}
\begin{tabularx}{\textwidth}{@{}>{\raggedright\arraybackslash}p{2cm}XX@{}}
\toprule
\textbf{Dimensão} & \textbf{Evidências documentais nas referências} & \textbf{Aplicação proposta ao Aula em Dia} \\
\midrule
Aquisição & Teachy e MagicSchool divulgam acesso inicial gratuito; Planbook apresenta período de avaliação de sua plataforma. & Página informativa objetiva, apresentação do benefício ao professor e caminho simples para cadastro. \\
Ativação & Teachy descreve a criação do primeiro planejamento; MagicSchool apresenta geração de plano a partir de objetivos; Planbook permite configurar turmas e começar a planejar. & Fluxo guiado para cadastrar disciplina e turma e salvar o primeiro plano de aula, com opção de gerar rascunho por IA. \\
Retenção & Teachy organiza aulas por calendário; Planbook permite copiar ou reaproveitar planejamentos; MagicSchool organiza ferramentas e coleções. & Histórico, edição e reutilização de planos, além de consulta por turma, disciplina e período. \\
Recomendação & Planbook documenta compartilhamento de planos; MagicSchool possibilita compartilhar coleções de ferramentas; Teachy anuncia colaboração no planejamento. & Compartilhamento voluntário de planos selecionados, com permissões e proteção de dados. \\
\bottomrule
\end{tabularx}
\par\smallskip
\begin{minipage}{\textwidth}\footnotesize Fonte: elaboração própria, fundamentada em \citet{teachy_planejamento2026,teachy_ferramentas2026,magicschool_ferramentas2026,magicschool_planos2026,planbook_professores2026}.\end{minipage}
\end{table}

\subsection{Indicadores propostos para validação futura}
Embora seja possível mapear oportunidades de produto sem acesso a dados privados das plataformas, a mensuração de AARR para o \textit{Aula em Dia} dependerá do desenvolvimento do sistema, da coleta legítima de informações e da autorização adequada dos participantes. Os indicadores a seguir são \textbf{propostas de avaliação}, ainda sem resultados reais:
\begin{itemize}
  \item \textbf{Aquisição:} número de participantes convidados que acessam a apresentação ou o protótipo e demonstram interesse em experimentar a plataforma.
  \item \textbf{Ativação:} proporção de participantes que concluem com sucesso o fluxo de criação e salvamento de seu primeiro plano de aula durante um teste autorizado.
  \item \textbf{Retenção:} proporção de participantes que, em uma avaliação longitudinal futura e consentida, retornam para consultar, editar ou criar novos planejamentos.
  \item \textbf{Recomendação:} proporção de participantes que, após testar a solução, declaram que a recomendariam a outro educador (medição por questão voluntária, sem coleta desnecessária de contatos de terceiros).
\end{itemize}
Os quatro indicadores acima não devem ser apresentados como métricas já coletadas. Os instrumentos de coleta, eventuais registros de uso e o tratamento de dados deverão respeitar as exigências da pesquisa e a legislação aplicável.

\subsection{Conclusões e oportunidades identificadas}
A análise documental sugere quatro diretrizes preliminares para o \textit{Aula em Dia}: (i) tornar a criação de planos simples e revisável pelo professor; (ii) centralizar planejamentos por turma e período, favorecendo a organização; (iii) viabilizar a reutilização e adaptação de planejamentos anteriores; e (iv) considerar compartilhamento opcional com controle de acesso. A referência à BNCC pode ser pertinente caso a pesquisa confirme que ela integra a rotina e as necessidades dos professores participantes.

Essas diretrizes são \textbf{hipóteses de requisitos}, a serem confrontadas com a pesquisa com usuários, Matriz CSD, histórias de usuário e protótipos. O benchmarking identifica possibilidades de projeto, mas não comprova aceitação, viabilidade de implementação nem superioridade em relação às plataformas comparadas.
